\chapter{نتیجه‌گیری و کارهای آینده}
\label{chap:conclusion}

در این فصل نخست در \autoref{sec:conclusion}، مروری بر آنچه در این پژوهش انجام شد صورت می‌پذیرد و نتایج اصلی جمع‌بندی می‌شوند. در ادامه نیز در \autoref{sec:futurework}، مسائل باز این حوزه و پیشنهاداتی برای کارهای آینده ارائه خواهد شد.

\section{نتیجه‌گیری}
\label{sec:conclusion}

این پژوهش به مسئله یادگیری فدرال در محیط‌های ابری چندگانه ناهمگون پرداخت. این محیط‌ها با سه چالش اساسی مشخص می‌شوند: ناهمگونی محاسباتی (تفاوت در توان پردازشی نمونه‌های ابری)، ناهمگونی شبکه (تفاوت در تأخیر و پهنای باند ارتباطات بین‌ناحیه‌ای)، و ناهمگونی داده (توزیع \gls{NonIID} داده‌های محلی). علاوه بر این، در محیط‌های ابری واقعی احتمال وجود گره‌های معیوب یا مخرب نیز وجود دارد.

روش‌های موجود یا تنها یکی از این چالش‌ها را هدف می‌گیرند، یا برای حل همزمان آن‌ها هزینه محاسباتی سنگینی دارند. در این پژوهش، مکانیزم \gls{TWAC} پیشنهاد شد که با طراحی یکپارچه به همه این چالش‌ها پاسخ می‌دهد.

\subsection{خلاصه روش پیشنهادی}

\gls{TWAC} برای هر گره یک امتیاز اعتماد نگهداری می‌کند که در هر دور ارتباطی بر اساس دو مؤلفه به‌روز می‌شود. مؤلفه همسویی جهتی، شباهت کسینوسی به‌روزرسانی گره را با میانگین متحرک نمایی جهت تجمیع‌شده‌های گذشته اندازه می‌گیرد. مؤلفه بزرگی منطقی، نرم به‌روزرسانی را نسبت به میانه نرم‌های همه گره‌ها ارزیابی می‌کند. ترکیب این دو مؤلفه گره‌هایی را که جهت‌معکوس‌سازی انجام می‌دهند (از طریق مؤلفه اول) یا نرم‌های غیرعادی بزرگ دارند (از طریق مؤلفه دوم) شناسایی می‌کند. امتیازات اعتماد با میانگین متحرک نمایی به‌روز می‌شوند که موجب می‌شود تأثیر رفتار گذرا کمتر باشد و رفتار مستمر وزن بیشتری در ارزیابی داشته باشد.

برای مدیریت گره‌های کُند، زمان مهلت تطبیقی بر اساس میانه زمان‌های پاسخ تعیین می‌شود و گره‌هایی که پاسخ نمی‌دهند، امتیاز اعتمادشان با یک ضریب کوچک کاهش می‌یابد. این طراحی از محرومیت دائمی گره‌هایی که به دلیل مشکلات گذرا در یک دور پاسخ ندادند جلوگیری می‌کند.

\subsection{خلاصه نتایج آزمایشی}

آزمایش‌ها در هشت سناریوی متفاوت روی مجموعه داده \gls{CIFAR}-10 با بیست گره انجام شد. نتایج به‌طور خلاصه عبارتند از:

\begin{itemize}
    \tick در سناریوی پایه (بدون حمله)، \gls{TWAC} با دقت $81.3\%$ بهترین عملکرد را دارد که $1\%$ بیشتر از \lr{FedAvg} و $5.6\%$ بیشتر از \lr{Multi-Krum} است.
    
    \tick در برابر حمله جهت‌معکوس‌سازی با $30\%$ گره مخرب، \gls{TWAC} دقت $70.2\%$ دارد در مقابل $18.7\%$ برای \lr{FedAvg}. این فاصله ۵۱ درصدی چشمگیرترین نتیجه آزمایش‌هاست.
    
    \tick در سناریوی ترکیبی (ناهمگونی شدید + حمله + گره کُند) که واقعی‌ترین آزمایش است، \gls{TWAC} با دقت $64.8\%$ از \lr{FedAvg} با $39.2\%$ و \lr{Multi-Krum} با $55.1\%$ پیشی می‌گیرد.
    
    \tick پیچیدگی محاسباتی \gls{TWAC} برابر $O(nd)$ است که با \lr{FedAvg} برابر و از \lr{Multi-Krum} به‌مراتب کمتر است. در عمل، زمان اجرا حدود $1.1\times$ بیشتر از \lr{FedAvg} است.
\end{itemize}

\subsection{بحث در مورد محدودیت‌ها}

با وجود نتایج مثبت، چند محدودیت در این پژوهش وجود دارد که باید صادقانه بیان شوند. نخست، ارزیابی روی یک مجموعه داده واحد (\gls{CIFAR}-10) انجام شده است و تعمیم‌پذیری نتایج به سایر حوزه‌ها نیاز به بررسی بیشتر دارد. دوم، فرض کفایت امتیاز اعتماد برای تشخیص گره‌های مخرب ممکن است در برابر حملات پیچیده‌تری که به‌طور دقیق رفتار گره‌های صادق را تقلید می‌کنند (مثل حملات بک‌دور) کارساز نباشد. سوم، انتخاب مقادیر بهینه پارامترهایی مثل $\gamma$ و $\beta$ ممکن است در سناریوهای مختلف متفاوت باشد.


\section{کارهای آینده}
\label{sec:futurework}

بر اساس نتایج و محدودیت‌های پژوهش حاضر، جهت‌های زیر برای تحقیقات آینده پیشنهاد می‌شود:

\subsection{گسترش به مجموعه داده‌های بزرگ‌تر}

یک گام طبیعی بعدی، ارزیابی \gls{TWAC} روی مجموعه داده‌های بزرگ‌تر مثل \lr{ImageNet} یا \lr{FEMNIST} است. در \lr{FEMNIST} که یک مجموعه داده فدرال واقعی است، ناهمگونی داده ذاتاً وجود دارد و نیازی به شبیه‌سازی مصنوعی ندارد.

\subsection{مقاوم‌سازی در برابر حملات پیچیده‌تر}

حملات بک‌دور (\lr{Backdoor Attack}) \cite{bagdasaryan2020backdoor} نوعی از حملات هستند که در آن‌ها گره مخرب به‌گونه‌ای عمل می‌کند که مدل روی اکثر داده‌ها درست رفتار کند، اما روی یک زیرمجموعه خاص (مثلاً تصاویر با یک نشان خاص) رفتار مخرب داشته باشد. کشف اینکه آیا مکانیزم امتیاز اعتماد \gls{TWAC} می‌تواند این نوع حملات را نیز شناسایی کند، جهت جالبی برای پژوهش آینده است.

\subsection{انتخاب تطبیقی پارامترها}

در حال حاضر پارامترهای \gls{TWAC} (مثل $\gamma$، $\alpha$، $\beta$) به‌صورت ثابت تنظیم می‌شوند. یک پیشرفت احتمالی، طراحی مکانیزمی برای تنظیم خودکار این پارامترها بر اساس خصوصیات مشاهده‌شده در هر دور است. مثلاً $\gamma$ می‌تواند در دورهای ابتدایی کوچک‌تر باشد (تا امتیازات سریع‌تر همگرا شوند) و در دورهای بعدی بزرگ‌تر (تا پایداری بیشتری داشته باشند).

\subsection{ترکیب با فشرده‌سازی ارتباطی}

در محیط‌های ابری با پهنای باند محدود، فشرده‌سازی به‌روزرسانی‌های مدل می‌تواند هزینه ارتباطی را کاهش دهد. بررسی سازگاری \gls{TWAC} با روش‌های فشرده‌سازی مثل کوانتیزاسیون (\lr{Quantization}) یا درندگی (\lr{Sparsification}) و اینکه آیا فشرده‌سازی بر کیفیت امتیازات اعتماد تأثیر می‌گذارد، یک سؤال تحقیقاتی جالب است.

\subsection{پیاده‌سازی توزیع‌شده واقعی}

پیاده‌سازی فعلی یک شبیه‌سازی روی یک ماشین است. یک گام مهم، پیاده‌سازی \gls{TWAC} در یک بستر توزیع‌شده واقعی مثل \lr{Flower} یا \lr{PySyft} و ارزیابی آن در یک محیط ابری چندگانه واقعی است. این پیاده‌سازی می‌تواند جنبه‌هایی از رفتار سیستم را آشکار کند که در شبیه‌سازی قابل مشاهده نیستند.

\subsection{تعمیم به یادگیری فدرال بدون سرور مرکزی}

در پیکربندی فعلی، یک سرور مرکزی برای تجمیع وجود دارد. یادگیری فدرال بدون سرور مرکزی (\lr{Decentralized FL}) که در آن گره‌ها مستقیماً با یکدیگر ارتباط برقرار می‌کنند، چالش جالبی برای تعمیم مفهوم امتیاز اعتماد ایجاد می‌کند. در این حالت، هر گره باید امتیاز اعتماد همسایگان خود را به‌صورت محلی نگهداری کند.

\begin{goal}{خلاصه کارهای آینده}
    کارهای پیشنهادی برای آینده را می‌توان در سه دسته کلی تقسیم‌بندی کرد: (الف) گسترش ارزیابی تجربی به مجموعه داده‌ها و سناریوهای حمله متنوع‌تر، (ب) بهبود الگوریتمی از طریق انتخاب تطبیقی پارامترها و ترکیب با فشرده‌سازی ارتباطی، و (ج) گسترش معماری به سمت سیستم‌های غیرمتمرکز و پیاده‌سازی واقعی در محیط‌های ابری.
\end{goal}